\documentclass{article}
\usepackage[utf8]{inputenc}
\usepackage[T1]{fontenc}
\usepackage[margin=2cm]{geometry}
\usepackage{hyperref}
\sloppy
\begin{document}
\section*{Historical source}
Encyclopaedia Britannica scan, source row 1942. Parse 3.5 extracted the text. Super 3.5 VL corrected it from an enlarged crop of each block. It can still contain recognition errors.
Page headers, footers, and page numbers are removed.

\subsection*{PHARMACY (E0170)}
PHARMACY.

Part II.

Prepara-
tions.

brought from Italy, are prepared in this manner. The
odorous parts may be again separated from the oil,
and transferred into water or fpirit, by distillation with
those liquors.

SECT. IV. Conservation of recent Vegetables and their Infufions, \&c. by Sugar and Honey.

§ 1. CONSERVES.

454. CONSERVES are compositions of recent vegetable matters and sugar beaten together into an uniform mass.

This management was introduced for preserving certain samples, undried, in an agreeable form, with as little alteration as possible in their native virtues; and to some subjects it is very advantageously applied. Vegetables, whose virtues are lost or destroyed in drying, may in this form be kept uninjured for a length of time: for, by carefully securing the mouth of the containing vessel, the alteration, as well as diffusion, of their active principles, is generally prevented; and the sugar preserves them from the corruption which juicy vegetables would otherwise undergo.

There are, however, sundry vegetables whole vi-
tous are impaired by this treatment. Muchaginous
fubilances, by long lying with sugar, become leis glu-
tionus, and that甚antly fully follow them the pa-
late. Many of the fingers flowers are of to tender
and delicate a texture, as almost entirely to lode their
peculiar qualities on being been on brushed
in general. It is obvious, that in this form, on ac-
count of the large admixture of sugar, only fulbliances
of considerable activity can be taken to advantage as
medicines. And indeed, conserve are at perfect con-
fidered chiefly as auxiliaries to medicines of greater
efficacy, as interventions for joining them together.
They are very convenient for echoing into boubles or
pills the more pondrous powders, as mecurunu dul-
cies, the colors of iron, and other mineral preparations;
which will liquid or leis confident matters, as fyrroup,
will not cohere.

The slops were formerly encumbered with many
conferences altogether insignificant; the few now retain-
ed have in general either an agreeable flavour to re-
commend them, or are capable of answering some
useful purpose as medicine. Their common dole is
the bulk of a nutmeg, or as much as can be taken up
at once or twice upon the point of a knife. There is
in general no great danger of exceeding in this parti-
cular.

455. General method of preferring conferences.

Leaves are picked from the stalks, and flowers from
their cups. They are beaten in a marble mortar,
with a wooden pestle, into a smooth mass; after
which thrice their weight of double-refined sugar is
added by degrees, and the beating continued till
they are uniformly mixed.

The sugar should be pulverized by itself, and passed through a sieve, before it is mixed with the vegetable mass; otherwise it cannot easily be reduced to sufficient fineness, so as to be duly incorporated. Some vegetables are scarce reducible to the requisite fineness by beating in a mortar; such is orange-peel. This is most conveniently rasped or grated off from the fruit, then well mixed with the sugar, and the com-

pound fet by in a close vessel for some weeks ; after
which it may be beaten smooth with considerably less
labour than at first. This peel, and red rose-buds,
are commonly ground in a wooden mill made for that
purpose.

456. Conserve of the leaves of garden fourvy-grass. L.

This is the only form that feury-grass in futultance
can be kept in without the total line of its virtues.
The conserve retains the full suite and virtue of the
herb for a very conferable length of time, as a year
two, provided the vessel be made perfectly cleo
and ice in a cool place. It must be given in Ecurbotic
habits three or four times a-day, so efficient.

457. Conserve of the leaves of wood-sorrel. L.

This is a very elegant and grateful conserve; in taste it is lightly acidulous, with a peculiar flavour, which some resemble to that of green-tea. It is taken occasionally for quenching thirst, and cooling the mouth and fauces, in hot distempers. It may be usefully joined to the foregoing preparation, whose virtue it somewhat promotes, at the same time that it improves the taste.

458. Conserve of the leaves of spearmint. I.

The conserve of mint retains the taste and virtues
of the herb. It is given in weaknesses of the stomach
and retchings to vomit; and not unfrequently does
service in some cases of this kind, where the warmer
and more active preparations of mint would be less
proper.

459. Conserve of the leaves of rue. L.

This conserve is given from a dram to half an ounce
in crudities of the prime via, for promoting digestion,
and in hysteric disorders: it gently stimulates the fo-
lids, attenuates viscid juices, and excites the natural
secretions. Some have had a great opinion of it, ta-
ken in a morning, as a preservative against the effects
of contagious air or exhalation.

460. Conference of the tops of sea-wormwood. L.

The conserve of wormwood has been celebrated in
dropsies: Matthiolus relates, that several persons were
cured by it of that distemper, without the assistance
of any other medicine. Where the disorder indeed
proceeds from a simple laxity or flaccidity of the fo-
lids, the continued use of this medicine may be of some
service, as it appears to be a not inelegant mild cor-
roborant. It is directed to be given in the dose of
half an ounce, about three hours before meals.

461. Conserve of the buds of red roses. L

This is a very agreeable and useful conserve. A dram or two, diluted in warm milk, are frequently given as a light refrigerant in weaknesses of the stomach, and likewise in coughs and phthisical complaints. In the German Ephemerides, examples are related of very dangerous phthises cured by the continued use of this medicine: in one of these cases, 20 pounds of the conserve were taken in the space of a month; and in another, upwards of 30. Riverius mentions several other instances of this kind.

462. Conference of rosemary-flowers. L.

462. Conference of rosemary-flowers. L.
Rosemary-flowers in great measure lose their peculiar


\section*{Source record}
\begin{description}
\item[Scan ID] \texttt{britannica:309f765b1ecc83e0a7f7762e44b4a927a7555443:train:1942}
\item[Image file] \texttt{page-1942-99b5dc3b4c6c.jpg}
\item[Image SHA-256] \texttt{99b5dc3b4c6cd19e10ef2da5ff4549329115fd8dba02d32d04fb855d1c774164}
\item[Dataset] NationalLibraryOfScotland/encyclopaedia\_britannica\_illustrated at revision \texttt{309f765b1ecc83e0a7f7762e44b4a927a7555443} (CC0-1.0)
\item[Extraction] model \texttt{nvidia/nemotron-3.5-parse}, raw reply \texttt{extract/raw/page-1942-99b5dc3b4c6c.json}
\item[Raw reply SHA-256] \texttt{2c1f89af03d97fd8969ac8fec7edc85691a052d127a526bd3afc61da74f28381}
\item[Correction] model \texttt{nvidia/Nemotron-3.5-Super-VL}, block crops, raw replies \texttt{correct/raw/crop\_nothink/page-1942-99b5dc3b4c6c-b*.json}; parse\_fallback: 1, super\_vl: 27
\item[Running header] none
\end{description}
\end{document}
